\section{Background and Problem Statement}\label{sec:background}

\subsection{UI Component Behavior as Relational Contracts}\label{sec:background:components}
UI component libraries expose compact declarative APIs, but the behavior behind those APIs is often rich and stateful~\cite{jartarghar2022react, lu2025misty, liu2024enhancing}. A component configured through a small set of props, events, or slots may involve rendering logic, state-dependent branches, accessibility attributes, layout calculations, and callback dispatch~\cite{zhang2024llamatouch}. Consequently, correctness often depends on interactions among component properties, states, and execution contexts rather than on individual executions alone.
Many UI component behaviors are therefore inherently relational. Correctness is determined not only by the outcome of a single execution, but also by how observable outputs change under variations in inputs, state, interaction, or context. For example, changing a placement option should consistently update the rendered position, while keyboard navigation should preserve focus and accessibility semantics. Assessing UI component correctness therefore requires reasoning about behavioral relations across executions rather than validating each execution in isolation.

\subsection{Execution Coverage and the Oracle Gap}\label{sec:background:coverage}
UI component libraries are commonly tested using rendering frameworks, interaction utilities, and coverage instrumentation tools. 
Existing UI component test suites commonly exercise behaviors through rendering, state updates, and simulated user interactions. However, exercising a behavior does not necessarily imply that the expected behavioral relation is explicitly validated.

Execution-based metrics, such as statement and branch coverage, indicate whether implementation code is exercised during testing, but provide limited evidence about whether the exercised behavior is behaviorally validated. A test that merely executes a branch and a test that explicitly verifies the behavioral relation implemented by that branch contribute equally to structural coverage. This limitation gives rise to an oracle gap: behavioral relations may be exercised without being explicitly checked. Consequently, high execution coverage alone does not necessarily indicate strong behavioral validation.

% UI component libraries are commonly tested using rendering frameworks, interaction utilities, and coverage instrumentation tools. Existing test suites often exercise component behavior through rendering, state updates, and user interactions. However, exercising a behavior does not necessarily mean that the expected behavioral property is explicitly verified.
% Structural coverage metrics, such as statement and branch coverage, measure whether implementation paths are executed during testing. While useful for identifying untested code regions, they do not indicate whether the exercised behavior is adequately validated. A test that merely reaches a branch and a test that verifies the behavioral relation implemented by that branch may contribute equally to structural coverage. This limitation gives rise to an oracle gap: component behaviors may be executed during testing without the corresponding behavioral relations being checked. Consequently, high execution coverage does not necessarily imply strong behavioral validation.

\subsection{Metamorphic Relations as an Adequacy Basis}\label{sec:background:mt}

The oracle gap arises because UI component correctness is often defined by relationships among executions rather than by the outcome of a single execution. Metamorphic relations (MRs) naturally represent such expectations by specifying how observable outputs should relate across executions whose inputs or execution contexts are systematically transformed.
For UI components, MRs capture behavioral expectations such as consistency, monotonicity, invariance, and accessibility-preserving transformations. Although these relations are often implicit in component specifications and developer expectations, they provide a useful behavioral reference for assessing existing test suites.

Rather than asking only which implementation code is executed, we can ask which inferred behavioral relations are exercised and explicitly validated. This motivates the central question of this work: can inferred component-level MRs serve as an empirical behavioral reference for assessing the behavioral adequacy of existing UI component test suites?